Compiled Stateflow gives the multi-GPU out-of-core trainer a
representation for decisions it previously hid: whether an admit is
sourced from host memory, kept in place, or handed across lanes. On
Twitter 16p, the compiler side is already decisive:
\textsc{Lane\_Matched} lowers selection cost by 14\%, cuts host admits
from 72 to 62, and replaces 10 \texttt{FULL\_RELOAD} events with 18
\texttt{PEER\_RELAY} descriptors.

The runtime side is now correct end-to-end. On two \SI{24}{GB} RTX
A5000s, the embeddings-scope peer runtime executes every planned relay
for both \textsc{Lane\_Matched} and \textsc{Disjoint\_Rounds} with zero
host fallback and zero descriptor mismatch, and
\textsc{Lane\_Matched} remains the strongest peer-enabled schedule.
What remains open is performance, not semantics: on this PCIe system,
peer execution displaces host traffic but leaves epoch time effectively
flat relative to the host-reload baseline.

That is still a useful checkpoint. The IR, compiler, validator, and
runtime contract are now stable enough to support the next phase:
bringing peer traffic under optimization with a peer-aware cost model
and reducing swap overhead through more overlapped execution.
